\begin{explanationbox}
	\textbf{\textcolor{CExplanation}{一句话直觉}}

	琴生不等式揭示了凸函数下的平均关系：加权平均的函数值不超过函数值的加权平均，等号统一意味着输入完全一致。

	\medskip
	\textbf{\textcolor{CExplanation}{核心拆解}}
	\begin{description}[style=nextline, leftmargin=2.8em, labelsep=0.5em, itemsep=0.45em]
		\item[\textbf{要点一（凸性定义）：}] 凸函数$f$满足$f(\lambda x+(1-\lambda)y)\le\lambda f(x)+(1-\lambda)f(y)$，等号当且仅当$x=y$。
		\item[\textbf{要点二（权重归一）：}] 权重$\lambda_i$之和为1，使$\sum\lambda_i x_i$成为$x_i$的加权平均。
		\item[\textbf{要点三（取等条件）：}] 严格凸时，等号成立$\iff$所有$x_i$相等；严格凹时方向相反，条件不变。
	\end{description}

	\medskip
	\textbf{\textcolor{CExplanation}{几何本质（函数凸性）}}

	凸函数图像位于弦的下方；弦两端点对应的函数连线，只有端点重合时二者重合。

	\medskip
	\textbf{\textcolor{CExplanation}{代数意义（加权平均关系）}}

	凸函数的加权平均输入的函数值，不大于函数值的加权平均；这反映了凸性在代换下的保序性。

	\medskip
	\textbf{\textcolor{CExplanation}{考点价值（不等式证明）}}
	\begin{itemize}[leftmargin=*, itemsep=0.5em, topsep=0.4em]
		\item \textbf{考法一（已知凹凸套用）：} 给定函数凸性，直接使用琴生不等式导出所需不等式。
		\item \textbf{考法二（构造凸函数）：} 将待证式变形为目标形式，选择合适的凸(凹)函数和权重。
	\end{itemize}

	\medskip
	\textbf{\textcolor{CExplanation}{顿悟点}}

	琴生不等式的灵魂是凸性；等号条件则是严格凸性不允许平直段的体现。

	\medskip
	\textbf{\textcolor{CExplanation}{使用场景}}
	\begin{itemize}[leftmargin=*, itemsep=0.5em, topsep=0.4em]
		\item \textbf{场景一（均值不等式）：} 取$f(x)=x^2$（凸），得平方平均$\ge$算术平均，等号当$x_i$相等。
		\item \textbf{场景二（指数/对数函数）：} $f(x)=e^x$（凸）或$f(x)=\ln x$（凹），可得到指数平均和对数不等式。
	\end{itemize}
\end{explanationbox}
